\section{Introduction}
Disaster response creates a demanding setting for aerial embodied intelligence. An inspection agent must interpret a task, place its sensor over a relevant region, decide whether the current evidence is sufficient, and return a usable assessment under limited observation and motion budgets. Static remote-sensing benchmarks measure recognition on fixed images, but they do not expose the decision of where and when to observe again. General aerial navigation environments provide motion and language tasks, yet they do not necessarily connect a changed view to a disaster-specific answer. A reproducible bridge between these settings is needed to study whether an agent's sensing actions improve the decision that motivated them.

A central challenge is how to transform static remote-sensing records into an environment in which an agent can act, reobserve, and be evaluated. Geographic metadata alone do not provide that environment. Disaster imagery, annotated buildings, task regions, vehicle states, and sensing actions must remain consistent within a shared world representation. This is particularly demanding for orthographic imagery: an apparent altitude change must respect the source-resolution ceiling, while the reference answer must remain tied to the task region even when the current view covers only part of it.

We present DisasterClaw, a georeferenced platform and task benchmark for studying task-dependent active observation in single-UAV disaster inspection. Its data-to-world pipeline places paired xBD records at their geographic footprints within a basemap and represents buildings, task ROIs, targets, and the UAV in the same coordinate system. Programmatic rules derive presence, damage, count, and spatial tasks from registered annotations and geometry. Reference answers are computed without a language model, and generated questions undergo automated checks and human review. The resulting world supports both interactive operation and headless benchmarking through shared observation, perception, action, and logging interfaces.

The distinction between a task and an observation is central to this design. Buildings are the units of damage perception, ROIs define the scope of questions, executed actions consume observation budgets, and complete episodes are scored by terminal answers. The agent can search, center on a target, descend to obtain a narrower footprint, answer, or abstain. Each executed view change is recorded together with the resulting evidence and answer. Generation variability is recorded separately because a vision-language model can change its answer even when environmental evidence is unchanged.

DisasterClaw is used here to examine budgeted reobservation. A closer view allocates more source pixels to the target region while discarding some surrounding context. Whether this helps depends on the task: a building label, a count, and a spatial relation require different evidence. We first characterize the ChangeOS binary perception response across three footprints on a development population. The main experiment compares seven frozen policies on a separate set of 160 questions reviewed independently by two annotators, scoring complete answers over three generation repeats.

Recent aerial embodied benchmarks emphasize broad skill coverage, realistic navigation, or standardized perception–decision–action evaluation, while recent remote-sensing agent benchmarks emphasize multimodal reasoning and tool use over large image collections. DisasterClaw targets a narrower complementary question: given a georeferenced disaster scene and a task-specific evidence requirement, can an executed observation action be traced to a change in perception evidence and ultimately to a corrected or harmed terminal answer? The contribution is therefore not broader environment realism or a new foundation model, but a controlled measurement interface for task-dependent observation value in disaster inspection.

The contributions are as follows.
\begin{enumerate}
\item We develop a map-anchored disaster-world construction pipeline that converts paired xBD imagery and building annotations into georeferenced interactive scenes. Task ROIs, vehicle states, targets, and movement actions share a geographic reference, enabling repeatable observation updates.
\item We construct an ROI-scoped disaster-inspection benchmark with four task families. Reference answers are derived deterministically from building labels and geometry. Automated schema, answer, geometry, duplicate, and split checks precede independent review by two annotators and adjudication.
\item We implement a traceable closed-loop inspection platform with an operator console and a headless runner. Vehicle states, observations, perception evidence, decisions, executed actions, budgets, and terminal answers connect each sensing decision to its observed task outcome.
\item We evaluate task-dependent active observation on 160 verified tasks and 3,360 policy episodes. T1\_TASK has the highest observed accuracy and the fewest reobservations among active policies in this sample, while adjusted intervals leave its advantage over holding and raw entropy uncertain. The results expose both benefits and harms of additional observations.
\end{enumerate}

The present platform is an imagery-based research environment, not a six-degree-of-freedom flight simulator. It uses orthographic source imagery and simplified motion, and it does not model wind, aircraft dynamics, occlusion, collision avoidance, or battery discharge. These boundaries determine the scope of every experimental claim.
